% ===========================================================================
% §4 — Architecture du calcul : notation, bases internes, ordre des postes
% ===========================================================================
\section{L'architecture du calcul}
\label{sec:architecture}

\subsection{Notation}

Les règles des sections suivantes utilisent une notation commune :

\begin{table}[ht]
\centering
\small
\begin{tabular}{c l}
\toprule
Symbole & Signification \\
\midrule
$r_i$            & revenu brut (travail ou retraite) de l'adulte $i \in \{1, 2\}$ \\
$R = r_1 + r_2$  & revenu brut du ménage ($r_2 = 0$ si un seul adulte) \\
$a_i$            & âge de l'adulte $i$ \\
$n$              & nombre d'enfants à charge \\
$\pos{x}$        & partie positive : $\pos{x} = \max(0,\,x)$ \\
$\RFN$           & revenu familial net (base québécoise, ligne 275) \\
$\AFNI$          & revenu familial net rajusté (base fédérale) \\
$\RAL$           & revenu aux fins de l'allocation-logement \\
\bottomrule
\end{tabular}
\caption{Notation utilisée dans les règles de calcul.}
\label{tab:notation}
\end{table}

Sauf indication contraire, les montants sont annuels, les calculs par adulte
sont ensuite additionnés sur le ménage, et le résultat final de chaque poste
est arrondi au cent.

\subsection{Les trois bases de revenu intermédiaires}
\label{sec:bases}

La subtilité centrale du système : les programmes ne sont pas modulés sur le
revenu brut, mais sur des bases de revenu \emph{construites par le calcul
lui-même}. Le moteur en reconstruit trois.

\paragraph{Revenu familial net (Québec).} Somme des revenus nets des adultes
au sens de la ligne 275 de la déclaration québécoise :
\begin{equation}
\label{eq:rfn}
\RFN \;=\; R \;+\; PSV \;+\; AS \;-\; D_{trav} \;-\; RRQ_{suppl}
\end{equation}
où $PSV$ est l'ensemble des prestations de la Sécurité de la vieillesse
(pension \emph{et} suppléments non imposables --- voir l'encadré ci-dessous),
$AS$ l'aide de dernier recours, $D_{trav}$ la déduction pour travailleur
(\pc{6} du revenu de travail, plafonnée --- poste~19) et $RRQ_{suppl}$ les
cotisations RRQ supplémentaires, déductibles. Cette base sert à la prime
RAMQ, à l'Allocation famille, à la prime au travail, au crédit pour la
solidarité, au soutien aux aînés et aux crédits modulés de l'impôt québécois.

\paragraph{Revenu familial net rajusté (fédéral, AFNI).} Même logique, sans
la déduction pour travailleur (propre au Québec), et diminuée de la déduction
fédérale pour frais de garde :
\begin{equation}
\label{eq:afni}
\AFNI \;=\; R \;+\; PSV \;+\; AS \;-\; RRQ_{suppl} \;-\; D_{garde}
\end{equation}
Cette base module l'Allocation canadienne pour enfants, le crédit TPS,
l'Allocation canadienne pour les travailleurs et le supplément médical
fédéral.

\paragraph{Revenu aux fins de l'allocation-logement.} Le programme
Allocation-logement applique aux aînés un ajustement qui retranche une
fraction des pensions :
\begin{equation}
\label{eq:ral}
\RAL \;=\; \RFN \;-\; PSV \times f
\end{equation}
où $f$ est un facteur d'ajustement propre à l'année et à la composition
(personne seule ou couple), extrait du calculateur officiel
(2025 : $f \approx \num{0.0496}$ ; 2026 : $f \approx \num{0.0389}$).

\paragraph{Encadré --- les prestations non imposables comptent quand même.}
Le Supplément de revenu garanti (SRG) et l'Allocation au conjoint ne sont pas
\emph{imposables}, mais ils doivent être \emph{inclus dans le revenu net}
servant à calculer les crédits d'impôt réductibles selon le revenu (montant
pour conjoint, montant en raison de l'âge, crédit médical) et la prime du
régime public d'assurance médicaments\footnote{Ministère des Finances du
Québec, \emph{Dépenses fiscales}, fiche 110113 :
\url{https://www.budget.finances.gouv.qc.ca/budget/outils/depenses-fiscales/fiches/fiche-110113.asp}}.
Le moteur distingue donc soigneusement, pour chaque adulte, le revenu
\emph{imposable} (sans SRG ni Allocation) du revenu \emph{net} (avec).

\subsection{L'ordre des postes}

Les dépendances imposent un ordre de calcul précis :

\begin{enumerate}
  \item \textbf{Cotisations sur le revenu brut} (postes 1--4) : RRQ, RQAP,
        assurance-emploi (ménages actifs), FSS (retraités). Elles ne
        dépendent que du revenu saisi ;
  \item \textbf{Postes-feuilles} : Sécurité de la vieillesse (poste~17,
        fonction du revenu de retraite et de l'âge) et aide de dernier
        recours (poste~13, fonction du revenu net des cotisations). L'aide
        sociale a besoin des cotisations ; c'est pourquoi elle vient après ;
  \item \textbf{Bases intermédiaires} : $\RFN$, $\AFNI$ et $\RAL$,
        équations~\eqref{eq:rfn} à~\eqref{eq:ral} ;
  \item \textbf{Prime RAMQ} (poste~5) sur le $\RFN$, en tenant compte des
        exonérations individuelles (aide sociale ; SRG quasi maximal) ;
  \item \textbf{Impôts} (poste~19) : le calcul québécois et le calcul
        fédéral consomment les cotisations (crédits et déductions), la PSV
        (revenu imposable et revenu net) et la prime RAMQ (crédit médical
        des couples) ;
  \item \textbf{Transferts} : tous les programmes modulés sur $\RFN$, $\AFNI$
        ou $\RAL$ (postes 6--11, 14--16) ;
  \item \textbf{Agrégation} (poste~20) : équation~\eqref{eq:rd}.
\end{enumerate}

Deux \og boucles \fg{} apparentes méritent explication. D'abord, l'aide
sociale est un \emph{intrant} d'autres programmes : elle entre dans $\RFN$ et
$\AFNI$ (elle est un revenu au sens de ces bases) et déclenche l'exonération
de la prime RAMQ. Ensuite, l'Allocation canadienne pour les travailleurs
dépend de l'$\AFNI$, qui inclut l'aide sociale : un dollar d'aide sociale
peut donc réduire l'ACT de vingt cents. Le moteur résout ces dépendances par
le simple ordre ci-dessus --- aucune itération n'est nécessaire, car aucun
programme ne dépend de sa propre sortie.

\subsection{Conventions de calcul}

\begin{itemize}
  \item \textbf{Annualisation} : les programmes versés mensuellement (aide
        sociale, allocation-logement) ou trimestriellement (Sécurité de la
        vieillesse) sont exprimés en montants annuels ; pour la Sécurité de
        la vieillesse, les paramètres annuels sont la moyenne des quatre
        barèmes trimestriels, d'où leurs décimales ;
  \item \textbf{Années de versement} : pour les prestations versées de
        juillet $Y$ à juin $Y{+}1$ (ACE, crédit TPS, solidarité), l'année
        affichée $Y$ correspond à cette période de versement ;
  \item \textbf{Arrondis} : chaque poste arrondit son résultat au cent,
        comme le calculateur officiel ; les règles internes d'arrondi
        particulières (tranches de 24/48/96~\$ du SRG, effort logement à
        quatre décimales) sont reproduites à l'identique ;
  \item \textbf{Cotisations en positif} : le moteur renvoie les cotisations
        et impôts en montants positifs ; c'est l'équation~\eqref{eq:rd} qui
        porte les signes.
\end{itemize}

% ============================================================================
% GÉNÉRÉ par scripts/exemples-guide.ts — NE PAS ÉDITER À LA MAIN.
% Régénérer : npm run exemples (chaque chiffre est vérifié contre le moteur).
% ============================================================================
\subsection{Les ménages types des exemples}
\label{sec:menages-exemples}

Les sections~\ref{sec:cotisations} à~\ref{sec:rd} illustrent chaque poste par
des exemples chiffrés, calculés pour l'\textbf{année d'imposition 2025} sur un
même jeu de treize ménages types (tableau~\ref{tab:menages-exemples}). Chaque
exemple reproduit l'algorithme du poste à la calculatrice ; tous les montants
sont produits par le moteur de calcul et chaque résultat final est
\textbf{vérifié automatiquement} contre celui-ci par la suite de tests du
projet (le fichier source des exemples est régénéré à chaque modification du
moteur).

\begin{table}[ht]
\centering
\small
\begin{tabular}{c l c l l}
\toprule
Code & Situation & Âges & Revenus bruts & Enfants (garde) \\
\midrule
M1 & Personne vivant seule & 25 & 0\,\$ & --- \\
M2 & Personne vivant seule & 30 & \D{9000} & --- \\
M3 & Personne vivant seule & 30 & \D{15000} & --- \\
M4 & Personne vivant seule & 40 & \D{50000} & --- \\
M5 & Personne vivant seule & 45 & \D{100000} & --- \\
M6 & Famille monoparentale & 35 & \D{35000} & 3 ans (subv., \D{2000}) \\
M7 & Couple & 45 / 44 & \D{30000} / 0\,\$ & --- \\
M8 & Couple & 38 / 36 & \D{60000} / \D{40000} & 4 ans (\D{13000}) ; 8 ans (\D{3000}) \\
M9 & Couple & 45 / 45 & \D{120000} / \D{60000} & 5 ans (\D{15000}) \\
M10 & Retraité vivant seul & 70 & \D{20000} & --- \\
M11 & Couple de retraités & 72 / 70 & \D{30000} / \D{10000} & --- \\
M12 & Retraité vivant seul & 76 & \D{110000} & --- \\
M13 & Couple de retraités & 68 / 66 & \D{12000} / \D{6000} & --- \\
\bottomrule
\end{tabular}
\caption{Les treize ménages types des exemples. Les revenus sont des revenus
de travail (ménages actifs) ou de pension privée (ménages retraités) ; les
frais de garde sont les frais annuels payés.}
\label{tab:menages-exemples}
\end{table}

Plusieurs postes sont modulés sur les bases de revenu de la
section~\ref{sec:bases} plutôt que sur le revenu brut. Le
tableau~\ref{tab:bases-exemples} donne ces bases, telles que le moteur les
reconstruit (équations~\eqref{eq:rfn} à~\eqref{eq:ral}) : les exemples y
renvoient sans refaire ce calcul, détaillé au poste~19 et à la
section~\ref{sec:bases}.

\begin{table}[ht]
\centering
\small
\begin{tabular}{c r r r r}
\toprule
Code & Revenu brut $R$ & $\RFN$ & $\AFNI$ & $\RAL$ \\
\midrule
M1 & \num{0} & \num{10548} & \num{10548} & \num{10548} \\
M2 & \num{9000} & \num{14388.77} & \num{14928.77} & \num{14388.77} \\
M3 & \num{15000} & \num{13985} & \num{14885} & \num{13985} \\
M4 & \num{50000} & \num{48115} & \num{49535} & \num{48115} \\
M5 & \num{100000} & \num{97506} & \num{98926} & \num{97506} \\
M6 & \num{35000} & \num{33265} & \num{32685} & \num{33265} \\
M7 & \num{30000} & \num{28315} & \num{29735} & \num{28315} \\
M8 & \num{100000} & \num{96230} & \num{86070} & \num{96230} \\
M9 & \num{180000} & \num{175521} & \num{170361} & \num{175521} \\
M10 & \num{20000} & \num{29891.99} & \num{29891.99} & \num{29401.05} \\
M11 & \num{40000} & \num{57582.28} & \num{57582.28} & \num{56708.72} \\
M12 & \num{110000} & \num{115737.85} & \num{115737.85} & \num{115453.08} \\
M13 & \num{18000} & \num{41237.58} & \num{41237.58} & \num{40083.04} \\
\bottomrule
\end{tabular}
\caption{Bases de revenu 2025 des ménages types, reconstruites par le moteur
(montants en dollars). Pour les retraités, $\RFN$ et $\AFNI$ incluent la
Sécurité de la vieillesse (pension, SRG et suppléments) ; $\RAL$ en retranche
une fraction (équation~\eqref{eq:ral}).}
\label{tab:bases-exemples}
\end{table}

